Model cards serve to disclose information about a trained machine learning model. This includes how it was built, what assumptions were made during its development, what type of model behavior different cultural, demographic, or phenotypic population groups may experience, and an evaluation of how well the model performs with respect to those groups. Here, we propose a set of sections that a model card should have, and details that can inform the stakeholders discussed in Section \ref{sec:motivation}.  A summary of all suggested sections is provided in Figure \ref{fig:model_sections}.

The proposed set of sections below are intended to provide relevant details to consider, but are not intended to be complete or exhaustive, and may be tailored depending on the model, context, and stakeholders.  Additional details may include, for example, interpretability approaches, such as saliency maps, TCAV \cite{KimEtAl2018}, and Path-Integrated Gradients \cite{PathIntegratedGradients2017,PathIntegratedGradients2018}); stakeholder-relevant explanations (e.g., informed by a careful consideration of philosophical, psychological, and other factors concerning what is as a good explanation in different contexts \cite{GoogleResponsiblePractices}); and privacy approaches used in model training and serving.

\begin{figure}
\begin{framed}
{\Large {\bf Model Card}}
\begin{itemize}[leftmargin=*]
\item {\bf Model Details}.  Basic information about the model.
\begin{itemize}
\item Person or organization developing model
\item Model date
\item Model version
\item Model type
\item Information about training algorithms, parameters, fairness constraints or other applied approaches, and features
\item Paper or other resource for more information
\item Citation details
\item License
\item Where to send questions or comments about the model \end{itemize}
\item {\bf Intended Use}.  Use cases that were envisioned during development. 
\begin{itemize}
\item Primary intended uses
\item Primary intended users
\item Out-of-scope use cases
\end{itemize}
\item {\bf Factors}.  Factors could include demographic or phenotypic groups, environmental conditions, technical attributes, or others listed in Section \ref{sec:factors}.  
\begin{itemize}
\item Relevant factors
\item Evaluation factors
\end{itemize}
\item {\bf Metrics}.  Metrics should be chosen to reflect potential real-world impacts of the model.
\begin{itemize}
\item Model performance measures
\item Decision thresholds
\item Variation approaches
\end{itemize}
\item {\bf Evaluation Data}.  Details on the dataset(s) used for the quantitative analyses in the card.
\begin{itemize}
\item Datasets
\item Motivation
\item Preprocessing
\end{itemize}
\item {\bf Training Data}.  May not be possible to provide in practice.  When possible, this section should mirror Evaluation Data.  If such detail is not possible, minimal allowable information should be provided here, such as details of the distribution over various factors in the training datasets.
\item {\bf Quantitative Analyses}
\begin{itemize}
\item Unitary results
\item Intersectional results
\end{itemize}
\item {\bf Ethical Considerations}
\item {\bf Caveats and Recommendations}
\vspace{-.25em}
\end{itemize}
\end{framed}
\vspace{-1em}
\caption{Summary of model card sections and suggested prompts for each.}\label{fig:model_sections}
\vspace{-1em}
\end{figure}
\subsection{Model Details}
This section of the model card should serve to answer basic questions regarding the model version, type and other details.  \\
{\bf Person or organization developing model}:  What person or organization developed the model? This can be used by all stakeholders to infer details pertaining to model development and potential conflicts of interest. \\
{\bf Model date}:  When was the model developed?  This is useful for all stakeholders to become further informed on what techniques and data sources were likely to be available during model development. \\
{\bf Model version}:  Which version of the model is it, and how does it differ from previous versions?  This is useful for all stakeholders to track whether the model is the latest version, associate known bugs to the correct model versions, and aid in model comparisons. \\
{\bf Model type}:   What type of model is it?  This includes basic model architecture details, such as whether it is a  Naive Bayes classifier, a Convolutional Neural Network, etc.  This is likely to be particularly relevant for software and model developers, as well as individuals knowledgeable about machine learning, to highlight what kinds of assumptions are encoded in the system.\\
{\bf Paper or other resource for more information}: Where can resources for more information be found?  \\
{\bf Citation details}: How should the model be cited? \\
{\bf License}: License information can be provided.  \\
{\bf Feedback on the model}:  E.g., what is an email address that people may write to for further information? \\

There are cases where some of this information may be sensitive. For example, the amount of detail corporations choose to disclose might be different from academic research groups. This section should not be seen as a requirement to compromise private information or reveal proprietary training techniques; rather, a place to disclose basic decisions and facts about the model that the organization can share with the broader community in order to better inform on what the model represents.

\subsection{Intended Use}
This section should allow readers to quickly grasp what the model should and should not be used for, and why it was created. 
It can also help frame the statistical analysis presented in the rest of the card, including a short description of the user(s), use-case(s), and context(s) for which the model was originally developed. Possible information includes: \\
{\bf Primary intended uses}: This section details whether the model was developed with general or specific tasks in mind (e.g., plant recognition worldwide or in the Pacific Northwest). The use cases may be as broadly or narrowly defined as the developers intend. For example, if the model was built simply to label images, then this task should be indicated as the primary intended use case. \\
{\bf Primary intended users}:  For example, was the model developed for entertainment purposes, for hobbyists, or enterprise solutions?  This helps users gain insight into how robust the model may be to different kinds of inputs. \\
{\bf Out-of-scope uses}: 
Here, the model card should highlight technology that the model might easily be confused with, or related contexts that users could try to apply the model to. This section may provide an opportunity to recommend a related or similar model that was designed to better meet that particular need, where possible. This section is inspired by warning labels on food and toys, and similar disclaimers presented in electronic datasheets. Examples include ``not for use on text examples shorter than 100 tokens'' or ``for use on black-and-white images only; please consider our research group's full-color-image classifier for color images.'' 


\subsection{Factors}\label{sec:factors}

Model cards ideally provide a summary of model performance across a variety of relevant factors including {\it groups}, {\it instrumentation}, and {\it environments}. We briefly describe each of these factors and their relevance followed by the corresponding prompts in the model card.

\subsubsection{Groups} ``Groups'' refers to distinct categories with similar characteristics that are present in the evaluation data instances. For human-centric machine learning models, ``groups'' are people who share one or multiple characteristics.  
Intersectional model analysis for human-centric models is inspired by the sociological concept of intersectionality, which explores how an individual's identity and experiences are shaped not just by unitary personal characteristics -- such as race, gender, sexual orientation or health -- but instead by a complex combination of many factors. These characteristics, which include but are not limited to cultural, demographic and phenotypic categories, are important to consider when evaluating machine learning models. Determining which groups to include in an intersectional analysis requires examining the intended use of the model and the context under which it may be deployed. Depending on the situation, certain groups may be more vulnerable than others to unjust or prejudicial treatment.

For human-centric computer vision models, the visual presentation of age, gender, and Fitzpatrick skin type~\cite{fitzpatrick} may be relevant. However, this must be balanced with the goal of preserving the privacy of individuals. As such, collaboration with policy, privacy, and legal experts is necessary in order to ascertain which groups may be responsibly inferred, and how that information should be stored and accessed (for example, using differential privacy \cite{DworkDifferential}).


Details pertaining to groups, including who annotated the training and evaluation datasets, instructions and compensation given to annotators, and inter-annotator agreement, should be provided as part of the data documentation made available with the dataset.  See \cite{GebruEtAl2018,HollandEtAl2018,statements} for more details.

\subsubsection{Instrumentation}
In addition to groups, the performance of a model can vary depending on what instruments were used to capture the input to the model. For example, a face detection model may perform differently depending on the camera's hardware and software, including lens, image stabilization, high dynamic range techniques, and background blurring for portrait mode. Performance may also vary across real or simulated traditional camera settings such as aperture, shutter speed and ISO. Similarly, video and audio input will be dependent on the choice of recording instruments and their settings. 
\subsubsection{Environment}
A further factor affecting model performance is the environment in which it is deployed.  For example, face detection systems are often less accurate under low lighting conditions or when the air is humid \cite{light}. Specifications across different lighting and moisture conditions would help users understand the impacts of these environmental factors on model performance.

\subsubsection{Card Prompts}
We propose that the Factors section of model cards expands on two prompts: \\
\noindent{\bf Relevant factors}:  What are foreseeable salient factors for which model performance may vary, and how were these determined? \\
{\bf Evaluation factors}:  Which factors are being reported, and why were these chosen?  If the relevant factors and evaluation factors are different, why? For example, while Fitzpatrick skin type is a relevant factor for face detection, an evaluation dataset annotated by skin type might not be available until reporting model performance across groups becomes standard practice.

\subsection{Metrics}\label{sec:metrics}
The appropriate metrics to feature in a model card depend on the type of model that is being tested. For example, classification systems in which the primary output is a class label differ significantly from systems whose primary output is a score.
In all cases, the reported metrics should be determined based on the model's structure and intended use. Details for this section include: \\
{\bf Model performance measures}:  What measures of model performance are being reported, and why were they selected over other measures of model performance?   \\
{\bf Decision thresholds}:  If decision thresholds are used, what are they, and why were those decision thresholds chosen?  When the model card is presented in a digital format, a threshold slider should ideally be available to view performance parameters across various decision thresholds.\\
{\bf Approaches to uncertainty and variability}:  How are the measurements and estimations of these metrics calculated? For example, this may include standard deviation, variance, confidence intervals, or KL divergence. Details of how these values are approximated should also be included (e.g., average of 5 runs, 10-fold cross-validation).
\subsubsection{Classification systems}
For classification systems, the error types that can be derived from a confusion matrix are {\it false positive rate}, {\it false negative rate}, {\it false discovery rate}, and {\it false omission rate}. We note that the relative importance of each of these metrics is system, product and context dependent. 

For example, in a surveillance scenario, surveillors may value a low false negative rate (or the rate at which the surveillance system fails to detect a person or an object when it should have). On the other hand, those being surveilled may value a low false positive rate (or the rate at which the surveillance system detects a person or an object when it should not have). 
We recommend listing all values and providing context about which were prioritized during development and why.

Equality between some of the different confusion matrix metrics is equivalent to some definitions of fairness. For example, equal false negative rates across groups is equivalent to fulfilling Equality of Opportunity, and equal false negative and false positive rates across groups is equivalent to fulfilling Equality of Odds \cite{HardtEtAl2016}.

\subsubsection{Score-based analyses}
For score-based systems such as pricing models and risk assessment algorithms, describing differences in the distribution of measured metrics across groups may be helpful. For example, reporting measures of central tendency such as the mode, median and mean, as well as measures of dispersion or variation such as the range, quartiles, absolute deviation, variance and standard deviation could facilitate the statistical commentary necessary to make more informed decisions about model development. A model card could even extend beyond these summary statistics to reveal other measures of differences between distributions such as cross entropy, perplexity, KL divergence and pinned area under the curve (pinned AUC)~\cite{DixonEtAl2018}.

There are a number of applications that do not appear to be score-based at first glance, but can be considered as such for the purposes of intersectional analysis. For instance, a model card for a translation system could compare BLEU scores~\cite{bleu} across demographic groups, and a model card for a speech recognition system could compare word-error rates. Although the primary outputs of these systems are not scores, looking at the score differences between populations may yield meaningful insights since comparing raw inputs quickly grows too complex. 

\subsubsection{Confidence}
Performance metrics that are disaggregated by various combinations of instrumentation, environments and groups makes it especially important to understand the confidence intervals for the reported metrics. Confidence intervals for metrics derived from confusion matrices can be calculated by treating the matrices as probabilistic models of system performance \cite{goutte2005probabilistic}.


\subsection{Evaluation Data}
All referenced datasets would ideally point to any set of documents that provide visibility into the source and composition of the dataset. Evaluation datasets should include datasets that are publicly available for third-party use. These could be existing datasets or new ones provided alongside the model card analyses to enable further benchmarking.  Potential details include: \\
{\bf Datasets}:  What datasets were used to evaluate the model?   \\
{\bf Motivation}:  Why were these datasets chosen?\\
{\bf Preprocessing}:  How was the data preprocessed for evaluation (e.g., tokenization of sentences, cropping of images, any filtering such as dropping images without faces)?  \\

To ensure that model cards are statistically accurate and verifiable, the evaluation datasets should not only be representative of the model's typical use cases but also anticipated test scenarios and challenging cases.  For instance, if a model is intended for use in a workplace that is phenotypically and demographically homogeneous, and trained on a dataset that is representative of the expected use case, it may be valuable to evaluate that model on two evaluation sets: one that matches the workplace's population, and another set that contains individuals that might be more challenging for the model (such as children, the elderly, and people from outside the typical workplace population). This methodology can highlight pathological issues that may not be evident in more routine testing. 

It is often difficult to find datasets that represent populations outside of the initial domain used in training. In some of these situations, synthetically generated datasets 
may provide representation for use cases that would otherwise go unevaluated~\cite{synthetic}. Section \ref{toxicity} provides an example of including synthetic data in the model evaluation dataset.

\subsection{Training Data}
Ideally, the model card would contain as much information about the training data as the evaluation data. However, there might be cases where it is not feasible to provide this level of detailed information about the training data. For example, the data may be proprietary, or require a non-disclosure agreement. In these cases, we advocate for basic details about the distributions over groups in the data, as well as any other details that could inform stakeholders on the kinds of biases the model may have encoded.
\subsection{Quantitative Analyses}
Quantitative analyses should be {\it disaggregated}, that is, broken down by the chosen factors.  Quantitative analyses should provide the results of evaluating the model according to the chosen metrics, providing confidence interval values when possible.  Parity on the different metrics across disaggregated population subgroups corresponds to how {\it fairness} is often defined \cite{verma2018fairness,MitchellEtAl2018}.  Quantitative analyses should demonstrate the metric variation (e.g., with error bars), as discussed in Section \ref{sec:metrics} and visualized in Figure \ref{fig:smiling_card}.  

\noindent The disaggregated evaluation includes: \\
{\bf Unitary results}:  How did the model perform with respect to each factor? \\
{\bf Intersectional results}:  How did the model perform with respect to the intersection of evaluated factors?

\subsection{Ethical Considerations}

This section is intended to demonstrate the ethical considerations that went into model development, surfacing ethical challenges and solutions to stakeholders.  Ethical analysis does not always lead to precise solutions, but the process of ethical contemplation is worthwhile to inform on responsible practices and next steps in future work.

While there are many frameworks for ethical decision-making in technology that can be adapted here \cite{markkula,deon,ethicalos}, the following are specific questions you may want to explore in this section: \\
{\bf Data}: Does the model use any sensitive data (e.g., protected classes)? \\
{\bf Human life}: Is the model intended to inform decisions about matters central to human life or flourishing -- e.g., health or safety?  Or could it be used in such a way? \\
{\bf Mitigations}:  What risk mitigation strategies were used during model development? \\
{\bf Risks and harms}: What risks may be present in model usage? Try to identify the potential recipients, likelihood, and magnitude of harms. If these cannot be determined, note that they were considered but remain unknown. \\
{\bf Use cases}: Are there any known model use cases that are especially fraught? This may connect directly to the intended use section of the model card.

If possible, this section should also include any additional ethical considerations that went into model development, for example, review by an external board, or testing with a specific community.

\subsection{Caveats and Recommendations}
This section should list additional concerns that were not covered in the previous sections. For example, did the results suggest any further testing?  Were there any relevant groups that were not represented in the evaluation dataset?  Are there additional recommendations for model use?  What are the ideal characteristics of an evaluation dataset for this model?